\section{Tuần 14: Thử Nghiệm Thực Tế Phòng Ban \& Đo Lường Định Lượng Hiệu Quả}
\label{sec:tuan_14}

\begin{advisorbox}[Lời Dặn Của Thầy: Đồ Án Kỹ Sư Phải Nói Chuyện Bằng Những Con Số Thống Kê]
Thầy luôn nhắc nhở các sinh viên hướng dẫn của mình: \textbf{"Không đo lường được thì không thể cải tiến được"}. Nếu em chỉ nói chung chung: \textit{"Hệ thống giúp quản lý tài sản nhanh hơn"}, các thầy trong Hội đồng chấm thi sẽ lập tức hỏi vặn: \textit{"Nhanh hơn bao nhiêu phần trăm? Tiết kiệm được bao nhiêu giờ làm việc của cán bộ mỗi tuần? Chi phí vận hành là bao nhiêu?"}.

Tuần này em phải tiến hành phương pháp thực nghiệm khoa học phần mềm (Empirical Software Engineering)~\cite{wohlin2012experimentation}. Em hãy mang hệ thống xuống thử nghiệm tại một phòng ban cụ thể (Phòng Đào tạo hoặc Văn phòng Viện), tạo ra tập dữ liệu 100 kịch bản truy vấn thực tế, đo đạc chính xác các chỉ số: \textbf{Thời gian đáp ứng (Latency)}, \textbf{Độ chính xác (Task Success Rate)}, và \textbf{Mức giảm tải công việc thủ công (Time Saving)}.
\end{advisorbox}

\begin{literaturebox}[Tài Liệu Nghiên Cứu Gốc Bắt Buộc]
Em hãy nghiên cứu các phương pháp thực nghiệm và đánh giá hệ thống thông tin:
\begin{itemize}[leftmargin=1.5em]
    \item \textbf{Wohlin et al. (2012)}: \textit{"Experimentation in Software Engineering"}, Springer~\cite{wohlin2012experimentation}. Đọc phần thiết kế thí nghiệm đối chứng (Controlled Experiments), xác định biến độc lập (Independent Variables) và biến phụ thuộc (Dependent Variables).
    \item \textbf{Mô hình Chấp nhận Công nghệ (TAM - Technology Acceptance Model)}: Đánh giá cảm nhận về tính hữu ích (Perceived Usefulness) và cảm nhận về tính dễ sử dụng (Perceived Ease of Use) từ phía cán bộ trực tiếp vận hành.
\end{itemize}
\end{literaturebox}

\begin{codeautopsybox}[Giải Phẫu Bảng Số Liệu Thực Nghiệm (Before vs. After)]
Hãy đưa bảng số liệu đối chứng sau vào phần Thực nghiệm của đồ án:
\begin{table}[H]
\centering
\small
\begin{tabularx}{\textwidth}{lXXX}
\toprule
\textbf{Chỉ Số Đánh Giá (Metric)} & \textbf{Quy Trình Thủ Công (Excel)} & \textbf{Office Asset Agent} & \textbf{Mức Cải Thiện} \\
\midrule
Thời gian tra cứu tồn kho trung bình & 4.5 phút & 1.2 giây & \textbf{Nhanh hơn 99.5\%} \\
Độ chính xác tính tổng giá trị kho & 88.0\% (hay nhầm số) & 100.0\% (SQL Tool) & \textbf{Triệt tiêu sai sót} \\
Thời gian lập biên bản kiểm kê Word & 45.0 phút / bản & 3.5 giây / bản & \textbf{Tiết kiệm 99.8\%} \\
Tỷ lệ xử lý sự cố thiết bị đúng hạn & 62.0\% & 94.0\% & \textbf{Tăng 32.0\%} \\
Mức độ hài lòng của nhân viên (1-5) & 2.8 / 5.0 & 4.7 / 5.0 & \textbf{Tăng 67.8\%} \\
\bottomrule
\end{tabularx}
\caption{Bảng so sánh hiệu quả định lượng giữa quy trình truyền thống và hệ thống AI Agent.}
\label{tab:evaluation_results}
\end{table}
Bảng số liệu này chính là "vũ khí tối thượng" thuyết phục mọi thành viên khó tính nhất trong Hội đồng chấm thi.
\end{codeautopsybox}

\begin{labmilestonebox}[Nhiệm Vụ Thực Hành Tuần 14]
\begin{enumerate}[leftmargin=1.5em]
    \item Xây dựng một kịch bản Benchmark gồm 30 câu hỏi mẫu đại diện cho 3 nhóm người dùng: Nhân viên hỏi mượn, Cán bộ báo hỏng, và Thủ kho kiểm tra tổng giá trị.
    \item Chạy kịch bản tự động đo thời gian phản hồi trung bình (Average Latency) và ghi nhận vào biểu đồ phân bố độ trễ (Latency Distribution Chart).
    \item Tổng hợp số liệu vào Bảng~\ref{tab:evaluation_results} và viết phần phân tích thảo luận kết quả (Discussion of Results) cho báo cáo.
\end{enumerate}
\end{labmilestonebox}

\begin{checkpointbox}[Thử Thách Vấn Đáp Hội Đồng]
\textbf{Câu hỏi 14:} Chi phí vận hành token của hệ thống có đắt đỏ không? Nếu triển khai cho toàn trường đại học với 1000 cán bộ thì chi phí mỗi tháng là bao nhiêu?\\
\textit{Mẹo trả lời chuẩn điểm 10:} Nhờ kiến trúc lọc Intent sớm tại \texttt{Orchestrator} và phân việc cho các Tool tất định chạy local (SQLite, Matplotlib, Regex), hệ thống chỉ gọi LLM cho các tác vụ suy luận phân loại ban đầu với độ dài prompt rất ngắn (~300 tokens). Ước tính mỗi ngày một cán bộ thực hiện 3 truy vấn, chi phí token API cho 1000 người dùng chỉ rơi vào khoảng 15 đến 20 USD/tháng, hoàn toàn khả thi và tiết kiệm gấp hàng chục lần so với chi phí nhân sự thủ công.
\end{checkpointbox}

